\section{Problem Statement and Motivation}

Vehicles in VANETs often encounter severe resource limitations when executing computational tasks locally \cite{wu2021design}. Modern vehicular applications—such as autonomous driving coordination, real-time hazard detection, cooperative traffic optimization, and sensor data processing—demand sophisticated processing capabilities that exceed the limited onboard computational resources, storage capacity, and battery power available in individual vehicles. This resource scarcity leads to increased task execution delays, higher failure rates for deadline-constrained tasks, and degraded quality of service for safety-critical applications. Consequently, task offloading to external computing infrastructure has become essential for meeting the stringent performance requirements of modern vehicular services.

To address these computational limitations, fog computing has emerged as a promising paradigm that deploys processing capabilities at Roadside Units (RSUs) and edge servers positioned throughout the network. By bringing computational resources closer to vehicles, fog computing reduces communication latency compared to traditional cloud-based solutions while providing sufficient processing power for complex vehicular applications. However, fog nodes themselves face resource constraints compared to cloud infrastructure, particularly during peak traffic periods in urban environments. This resource scarcity becomes critical when multiple vehicles repeatedly request identical computations. For instance, during rush hour at busy intersections, dozens of vehicles may independently request the same traffic pattern analysis, route optimization for common destinations, or environmental data processing from overlapping sensor coverage areas. Each redundant execution consumes valuable computational resources and energy at fog nodes, preventing these resources from serving other critical tasks and leading to increased delays, higher failure rates, and degraded system performance.

Task result caching emerges as a promising solution by storing outputs of previously executed tasks and reusing them for subsequent identical requests. By preventing repetitive execution, caching conserves both energy and computational resources at fog nodes, freeing these resources for executing other critical tasks and improving overall system throughput. When a vehicle submits a task request, the fog node first queries its cache; upon finding a matching result, it is immediately returned without re-execution, dramatically reducing response time, energy consumption, and computational overhead. However, implementing effective task result caching in VANETs introduces three critical challenges that existing solutions fail to address comprehensively.

First, determining which task results to cache proactively requires predicting future task popularity patterns across distributed fog nodes. Without intelligent pre-placement strategies, cache misses force redundant re-execution, wasting computational resources and energy. Traditional frequency-based heuristics cannot capture the complex spatiotemporal dynamics of vehicular environments where task popularity fluctuates rapidly based on geographic location, traffic conditions, and time of day. A task frequently requested in one area may become obsolete within seconds as vehicles move to different locations. Moreover, certain urban areas such as commercial centers, administrative districts, and recreational zones typically experience higher traffic density and generate more computational tasks compared to residential areas, creating uneven demand patterns that simple caching policies cannot handle effectively.

Second, limited storage capacity at fog nodes necessitates intelligent eviction policies to determine which cached results to retain as new tasks continuously arrive. Static policies such as LRU, LFU, and FIFO rely on simple rules that fail to account for task heterogeneity—including varying priorities, deadlines, computational costs, and data sizes. Naive eviction may discard valuable results that will be requested again shortly, leading to cache thrashing and degraded performance. The challenge intensifies in VANETs where vehicle mobility creates highly dynamic access patterns; a task frequently requested in one area may become obsolete within seconds as vehicles move to different locations. Real-time eviction decisions must balance multiple competing objectives: maximizing cache hit rates, minimizing task completion times, ensuring deadline compliance for urgent tasks, and conserving energy for overall system efficiency.

Third, coordinating caching strategies across distributed fog infrastructure requires centralized visibility and control that traditional distributed approaches cannot provide. The SDN controller maintains a global view of network conditions, task distributions, and resource availability across all RSUs, enabling system-wide optimization. However, leveraging SDN for intelligent, learning-based cache management in vehicular fog computing remains largely unexplored. Existing SDN-VANET solutions focus primarily on routing optimization and load balancing, overlooking the potential for coordinated caching decisions to amplify overall system performance. Furthermore, while machine learning techniques offer promising approaches for adaptive caching, existing research predominantly explores Federated Learning and Reinforcement Learning independently, failing to leverage their complementary strengths—FL for long-term pattern recognition and RL for short-term dynamic optimization.

In summary, effective task result caching in SDN-based VANETs requires a unified framework that addresses three key requirements: (1) proactive cache pre-placement based on predicted task popularity patterns, (2) adaptive real-time eviction policies that consider task heterogeneity and dynamic network conditions, and (3) coordinated decision-making across distributed fog nodes through centralized SDN control. Due to the lack of an existing study that addresses these requirements holistically, this paper proposes FL-RL-TRC; a novel two-phase intelligent caching architecture that synergistically combines Federated Learning for anticipatory cache population with Deep Reinforcement Learning for reactive eviction control, achieving superior cache performance while reducing energy consumption and improving overall network efficiency.
